\section{Conclusion}
\label{sec:conclusion}

% Paragraph 1: Main finding
A single randomly assigned tournament opportunity in professional tennis causes gains in tournament access but does not uniformly improve playing ability. In the full ATP Grand Slam sample ($N = 248$), a stacked specification shows that lottery selection increases ranking points by 31.9 at 4 weeks ($p = 0.028$), 38.5 at 12 weeks ($p = 0.007$), and 49.2 at 26 weeks ($p = 0.054$), with main draw entries increasing by 1.1 over 26 weeks ($p = 0.006$). On the WTA tour ($N = 132$), effects are larger at short horizons---64.4 points at 4 weeks ($p = 0.003$) and 78.7 at 8 weeks ($p < 0.001$)---but fade by 26 weeks. At the immediate event, treated players are 32 percentage points more likely to win a main draw match, earning an average of 29.66 ranking points. Elo ratings---which adjust for opponent quality---are unchanged at every horizon on both tours. A match-level tournament performance model finds no significant effect on match win probability on either tour: $\hat{\delta} = -0.123$ ($p = 0.119$; 8,543 matches from 130 players) on the ATP tour and $\hat{\delta} = +0.034$ ($p = 0.760$; 5,391 matches from 98 players) on the WTA tour. Ranking point gains come from additional participation rather than improved match performance. The full estimation sample (including repeat LL candidates) produces qualitatively similar results, with expected attenuation. The non-Grand-Slam control function estimates corroborate these findings: ATP ranking point gains of 12.2 at 4 weeks ($p = 0.065$) and 28.1 at 26 weeks ($p = 0.009$), with significant tournament access gains at longer horizons.

% Paragraph 2: Interpretation
The dominant mechanism is access accumulation. The ranking system does not adjust for opponent quality: a first-round win at a Grand Slam earns the same points regardless of the opponent's strength. When a lucky loser earns points at the focal event, their improved ranking grants entry to subsequent tournaments, generating further opportunities to accumulate points. Match-level performance conditional on opponent strength is unchanged on both tours: treated players perform neither better nor worse than controls once in the draw. Ranking point gains on both tours come from more draws entered and more matches played, not from improved competitive ability.

% Paragraph 3: Broader implications
These findings have implications for any labor market where threshold-based access generates cumulative advantage. Second-chance mechanisms---wildcard entries in sports, waitlist admissions in universities, trial hiring periods in firms---can provide real access benefits to marginal candidates. But when performance is continuously measured and positions are regularly updated, the signal from sustained ability dominates the noise from a single favorable shock. The contrast with \citet{Oyer2006_initial_placement} and \citet{Oreopoulos2012_graduating_recession}, where initial conditions cause persistent career effects in less transparent markets, suggests that persistence depends on how transparently performance is measured and aggregated.

% Paragraph 4: Limitations and future work
Several limitations apply. First, the lottery contamination caveat---that an unknown fraction of observations were assigned by ranking rather than by draw---biases toward attenuation but cannot be fully resolved without withdrawal timing data; the verified lottery subsample and first-LL-only restriction mitigate this concern. Second, the WTA Grand Slam sample ($N = 132$) has limited power, and WTA dynamic effects should be interpreted with caution. Third, the non-Grand-Slam control function estimates rely on the exact enumeration selection probability $P_i^{LL}$ for identification and may not generalize to all qualifying losers. Fourth, I measure tournament access and match performance, which are imperfect proxies for the career outcomes that matter most to players: earnings, career length, and sponsorship. Future work could examine longer-run outcomes---career length, total earnings, peak ranking achieved---to determine whether the temporary access boost has lasting consequences through sponsorship or confidence channels. The finding that access shocks generate ranking gains without affecting match-level ability deserves study in other settings with lottery-based or threshold-based access to test the generality of the result.
